\chapter{The Explicit Formula and the Riemann Hypothesis}

\begin{center}\itshape
If I were to awaken after having slept for a thousand years, my first question would be: has the Riemann hypothesis been proven?\\
\upshape --- attributed to David Hilbert
\end{center}

Near the end of his eight pages, Riemann wrote down a formula that nobody before him could have imagined. On the left side stands the prime counting
function, a staircase that climbs by one at every prime and is flat in between. On the right side stands a smooth main term, minus a sum of waves, one for each
zero of $\zeta$ in the critical strip. The staircase, with all its corners, \emph{is} the sum of the waves. Riemann did not prove it. Von Mangoldt did, in 1895,
for a slightly better-behaved staircase: the function $\psi$, which climbs by $\log p$ at every prime power $p^k$. In that form the formula reads
\[
\psi(x)=x-\sum_\rho\frac{x^\rho}\rho-\log2\pi-\frac12\log\Bigl(1-\frac1{x^2}\Bigr),
\]
the sum running over the non-trivial zeros $\rho$, taken in conjugate pairs. Each pair contributes $-2\operatorname{Re}\frac{x^\rho}\rho$. Writing $\rho=\beta+i\gamma$ and $x=e^u$, this is a wave of
amplitude $\frac{2e^{\beta u}}{|\rho|}$ and frequency $\gamma$ in the variable $u=\log x$. The zeros are the notes, and the primes are the music they make together.

Figure~\ref{fig:explicit} plays the music. With ten pairs of zeros the staircase is a wobbly ramp; with a hundred, the corners appear, exactly at the prime powers, with
steps of height exactly $\log p$. Nothing in the zeros seems to know about primes, and yet they draw them.

\begin{figure}[ht]
\centering
\fig{ch07-explicit}
\caption{The staircase $\psi(x)$ (blue) for $1.5\le x\le30$, and the right side of the explicit formula truncated to the first $10$ pairs of zeros (green, dashed) and to the
first $100$ pairs (orange). With $100$ pairs, the error is below $0.34$ at all points more than $0.3$ away from an integer.}
\label{fig:explicit}
\end{figure}

The formula also explains why the zeros matter so much. A zero with real part $\beta$ makes a wave of size about $x^\beta$; the prime number theorem says that
$\psi(x)-x$ is small compared with $x$, and the waves say how small, in terms of how far to the right the zeros go. This chapter proves the formula, then draws
three consequences: the error term of de la Vallée Poussin (1899), the equivalence between the Riemann hypothesis and the best possible error term (von Koch, 1901),
and the explanation, promised in Chapter~4, of why primes $\equiv3\pmod4$ win the race against primes $\equiv1\pmod4$ almost all the time.

The method is the one of Chapter~5 with sharper tools: write $\psi$ as a contour integral of $-\frac{\zeta'}\zeta$, push the contour to the left across the critical strip,
and collect a residue at every zero. Chapter~6 supplied what this needs: the partial fractions of $\frac{\zeta'}\zeta$, the count of the zeros, and a zero-free region.

\begin{think}
Look at the formula above at a prime $p$. The left side jumps by $\log p$; the terms $x$, $\log2\pi$ and $\frac12\log(1-x^{-2})$ are continuous. So the sum over zeros must jump.
But each term $\frac{x^\rho}\rho$ is a smooth function of $x$. What does this tell you about the convergence of $\sum_\rho\frac{x^\rho}\rho$?
\end{think}

\section{Perron's formula}

The first step turns a sum over integers into an integral. The tool is a single integral that acts as a switch: it equals $1$ when $y>1$ and $0$ when $y<1$.
For $y>0$ let $\delta(y)=1$ if $y>1$, $\delta(1)=\frac12$, and $\delta(y)=0$ if $y<1$. For $c>0$ and $T>0$ let
\[
I(y,T)=\frac1{2\pi i}\int_{c-iT}^{c+iT}\frac{y^s}s\dd s,
\]
where $y^s=e^{s\log y}$ and the path is the vertical segment.

\begin{lemma}[Perron]\label{lem:perron}
Let $y>0$, $c>0$, $T>0$. If $y\ne1$,
\[
\bigl|I(y,T)-\delta(y)\bigr|\le y^c\min\Bigl(1,\frac1{\pi T|\log y|}\Bigr).
\]
If $y=1$, $\bigl|I(1,T)-\frac12\bigr|\le\frac c{\pi T}$.
\end{lemma}
\begin{proof}
\emph{Case $y>1$.} Then $|y^s|=y^{\operatorname{Re}s}$ is small far to the left, so we push the segment to the left. For $U>0$, the function $\frac{y^s}s$ is holomorphic except for a simple pole
at $0$ with residue $1$. By the residue theorem (Tome~4, Theorem~9.26) on the rectangle with vertices $c\pm iT$ and $-U\pm iT$, $I(y,T)-1$ is $-\frac1{2\pi i}$ times the integral over
the three other sides, traversed from $c+iT$ to $-U+iT$ to $-U-iT$ to $c-iT$. On a horizontal side, $|s|\ge T$, so its integral is at most
$\frac1T\int_{-\infty}^cy^\sigma\dd\sigma=\frac{y^c}{T\log y}$. The left side has length $2T$ and the integrand is at most $\frac{y^{-U}}U$ there, which tends to $0$ as $U\to\infty$. So
$|I(y,T)-1|\le\frac1{2\pi}\cdot\frac{2y^c}{T\log y}=\frac{y^c}{\pi T\log y}$.

For the other bound, replace the rectangle by the arc of the circle $|s|=R$, $R=|c+iT|$, that goes from $c+iT$ to $c-iT$ through the left half-plane. The segment and the arc enclose $0$. On the
arc, $\operatorname{Re}s\le c$, so $\bigl|\frac{y^s}s\bigr|\le\frac{y^c}R$, and the arc has length less than $2\pi R$. So $|I(y,T)-1|\le\frac1{2\pi}\cdot2\pi R\cdot\frac{y^c}R=y^c$.

\emph{Case $y<1$.} Now $|y^s|$ is small far to the right. The same two contours, drawn to the right of the segment, enclose no pole, and the same estimates give
$|I(y,T)|\le\frac{y^c}{\pi T|\log y|}$ and $|I(y,T)|\le y^c$.

\emph{Case $y=1$.} Directly, with $s=c+it$,
\[
\begin{aligned}I(1,T)&=\frac1{2\pi}\int_{-T}^T\frac{\mathrm dt}{c+it}=\frac1{2\pi}\int_{-T}^T\frac{c\dd t}{c^2+t^2}\\&=\frac1\pi\arctan\frac Tc,\end{aligned}
\]
the odd part $\frac{-it}{c^2+t^2}$ integrating to $0$. So $\frac12-I(1,T)=\frac1\pi\bigl(\frac\pi2-\arctan\frac Tc\bigr)=\frac1\pi\arctan\frac cT\le\frac c{\pi T}$.
\end{proof}

The switch is not perfect: near $y=1$ it is blurred over a window of width about $\frac1T$ in $\log y$. That blur is the price of truncating at height $T$, and it is what the error terms
below pay for.

Since $\psi$ jumps at prime powers, the natural value there is the midpoint. Let
\[
\psi_0(x)=\psi(x)-\frac12\Lambda(x)
\]
when $x$ is an integer, and $\psi_0(x)=\psi(x)$ otherwise. For $x\ge2$, let $\langle x\rangle$ be the distance from $x$ to the nearest prime power \emph{other than} $x$ itself.

\begin{proposition}[truncated Perron formula]\label{prop:truncperron}
Let $x\ge2$, $T\ge2$, and $c=1+\frac1{\log x}$. Then
\[
\begin{aligned}\psi_0(x)&=\frac1{2\pi i}\int_{c-iT}^{c+iT}\Bigl(-\frac{\zeta'}\zeta(s)\Bigr)\frac{x^s}s\dd s\eqbr+O\Bigl(\frac{x\log^2x}T+\log x\min\Bigl(1,\frac x{T\langle x\rangle}\Bigr)\Bigr).\end{aligned}
\]
The implied constant is absolute.
\end{proposition}
\begin{proof}
On the line $\operatorname{Re}s=c$, $-\frac{\zeta'}\zeta(s)=\sum_n\Lambda(n)n^{-s}$ converges absolutely and uniformly (Lemma~\ref{lem:logderiv}), so we may integrate term by term:
the integral is $\sum_n\Lambda(n)I(\frac xn,T)$. Since $\psi_0(x)=\sum_n\Lambda(n)\delta(\frac xn)$, Lemma~\ref{lem:perron} bounds the difference $R$ by
\[
\begin{aligned}|R|&\le\sum_{n\ne x}\Lambda(n)\Bigl(\frac xn\Bigr)^c\min\Bigl(1,\frac1{\pi T|\log\frac xn|}\Bigr)\eqbr+\frac{c\Lambda(x)}{\pi T},\end{aligned}
\]
with $\Lambda(x)=0$ if $x$ is not an integer. Note $x^c=ex$ and $c\le1+\frac1{\log2}<3$. The last term is $O(\frac{\log x}T)$. Split the sum.

\emph{Far from $x$: $n\le\frac x2$ or $n\ge2x$.} Then $|\log\frac xn|\ge\log2$, and these terms contribute at most
\[
\frac{x^c}{\pi T\log2}\sum_n\frac{\Lambda(n)}{n^c}=\frac{ex}{\pi T\log2}\Bigl(-\frac{\zeta'}\zeta(c)\Bigr).
\]
By the second inequality in the proof of Theorem~\ref{thm:zerofree}, $-\frac{\zeta'}\zeta(c)\le\frac1{c-1}+A=\log x+A$ when $c\le2$, and it is bounded by $\sum_n\Lambda(n)n^{-2}$ when $c\ge2$. This part is $O(\frac{x\log x}T)$.

\emph{Near $x$: $\frac x2<n<2x$, $n\ne x$.} Here $(\frac xn)^c\le2^3$ and $\Lambda(n)\le\log2x$. Write $n=x(1+u)$ with $-\frac12<u<1$. For $u>0$, $\log(1+u)\ge\frac u{1+u}\ge\frac u2$; for $u<0$,
$|\log(1+u)|\ge|u|$. So $|\log\frac xn|\ge\frac{|n-x|}{2x}$. The integers $n$ with $|n-x|\ge1$ in this range contribute at most
\[
8\log(2x)\sum_{1\le|n-x|<2x}\frac{2x}{\pi T|n-x|}=O\Bigl(\frac{x\log^2x}T\Bigr),
\]
since for each $k\ge1$ at most two integers have $k\le|n-x|<k+1$. At most two integers $n\ne x$ have $|n-x|<1$. If such an $n$ is a prime power, then $|n-x|\ge\langle x\rangle$ and its term
is at most $8\log(2x)\min\bigl(1,\frac{2x}{\pi T\langle x\rangle}\bigr)$.
\end{proof}

\section{The explicit formula}

Now we move the segment to the left, all the way to $-\infty$. Three things must be controlled: the integrand on the far left, on the horizontal sides crossing the strip, and the residues.
The horizontal sides are the delicate part, since they cut through the critical strip where $\frac{\zeta'}\zeta$ has poles. The trick is to choose the height so that it passes
between the zeros, as far from all of them as their density allows.

\begin{lemma}[good heights]\label{lem:goodT}
There is an absolute constant $C$ such that for every $T\ge3$ there is $T'\in[T,T+1]$ with
\[
\Bigl|\frac{\zeta'}\zeta(\sigma\pm iT')\Bigr|\le C\log^2T\qquad(-1\le\sigma\le3).
\]
\end{lemma}
\begin{proof}
By Corollary~\ref{cor:nunit}, the number $m$ of ordinates in $[T,T+1]$ is at most $C_0\log T$. They cut $[T,T+1]$ into at most $m+1$ intervals, one of which has length at least $\frac1{m+1}$; let $T'$ be
its midpoint. Every ordinate $\gamma$ lies outside the open interval, so $|T'-\gamma|\ge\frac1{2(m+1)}$. For $-1\le\sigma\le2$, Proposition~\ref{prop:nearzeros} writes $\frac{\zeta'}\zeta(\sigma+iT')$ as
$\sum_{|\gamma-T'|\le1}\frac1{s-\rho}+O(\log T)$. The sum has $O(\log T)$ terms (Corollary~\ref{cor:nunit} again), each of modulus at most $\frac1{|T'-\gamma|}\le2(m+1)=O(\log T)$. For $2\le\sigma\le3$,
$|\frac{\zeta'}\zeta(s)|\le\sum_n\Lambda(n)n^{-2}$. Finally $\frac{\zeta'}\zeta(\overline s)=\overline{\frac{\zeta'}\zeta(s)}$ handles $-T'$.
\end{proof}

On the far left, the functional equation reflects $\frac{\zeta'}\zeta$ into the half-plane of absolute convergence, and only the Gamma factors remain.

\begin{lemma}[the left half-plane]\label{lem:leftzeta}
For $\operatorname{Re}s\le-1$ and $|s+2k|\ge1$ for every integer $k\ge1$,
\[
\frac{\zeta'}\zeta(s)=O\bigl(\log(|s|+2)\bigr).
\]
\end{lemma}
\begin{proof}
Take the logarithmic derivative of $\pi^{-s/2}\Gamma(\frac s2)\zeta(s)=\pi^{-(1-s)/2}\Gamma(\frac{1-s}2)\zeta(1-s)$ (Tome~4, Theorem~10.8):
\[
\begin{aligned}\frac{\zeta'}\zeta(s)&=\log\pi-\frac12\psi_\Gamma\Bigl(\frac s2\Bigr)\eqbr-\frac12\psi_\Gamma\Bigl(\frac{1-s}2\Bigr)-\frac{\zeta'}\zeta(1-s).\end{aligned}
\]
Here $\operatorname{Re}(1-s)\ge2$, so $|\frac{\zeta'}\zeta(1-s)|\le\sum_n\Lambda(n)n^{-2}$. The point $\frac{1-s}2$ has real part at least $1$, so Theorem~\ref{thm:stirlingc}(c) (with $\delta=\frac\pi2$) gives
$\psi_\Gamma(\frac{1-s}2)=\log\frac{1-s}2+O(1)=O(\log(|s|+2))$. For $\psi_\Gamma(\frac s2)$ we first reflect. Taking the logarithmic derivative of $\Gamma(z)\Gamma(1-z)=\frac\pi{\sin\pi z}$
(Tome~4, Theorem~5.13) gives $\psi_\Gamma(z)=\psi_\Gamma(1-z)-\pi\cot\pi z$. With $z=\frac s2$: $1-\frac s2$ has real part at least $\frac32$, so $\psi_\Gamma(1-\frac s2)=O(\log(|s|+2))$ as before. And $z$
stays at distance at least $\frac12$ from the integers $\le-1$, and at distance at least $\frac12$ from $0$ and the positive integers since $\operatorname{Re}z\le-\frac12$. The function $\cot\pi z$ is $1$-periodic,
continuous away from the integers, and tends to $\mp i$ as $\operatorname{Im}z\to\pm\infty$, so it is bounded on this set.
\end{proof}

Two constants will appear in the residues. First, $\zeta(0)=-\frac12\ne0$ (Tome~4, Chapter~10), and
\begin{equation}\label{eq:zetaprime0}
\frac{\zeta'}\zeta(0)=\log2\pi.
\end{equation}
Indeed, in \eqref{eq:zetapf} let $s\to0$. The sum over zeros vanishes at $s=0$, $-\frac1{s-1}\to1$, and $\psi_\Gamma(z)=\psi_\Gamma(z+1)-\frac1z=-\frac1z-\gamma+O(z)$ (from $\Gamma(z+1)=z\Gamma(z)$ and
$\psi_\Gamma(1)=-\gamma$), so $-\frac1s-\frac12\psi_\Gamma(\frac s2)\to\frac\gamma2$. Hence $\frac{\zeta'}\zeta(0)=B+1+\frac12\log\pi+\frac\gamma2$, and with $B=\frac12\log4\pi-1-\frac\gamma2$
(Exercise~\ref{ex:digammahalf}) this is $\frac12\log4\pi+\frac12\log\pi=\log2\pi$. Second, the trivial zeros $-2,-4,\dots$ are simple: they are the poles of $\Gamma(\frac s2)$, which are simple, at points
where $\xi\ne0$.

\begin{theorem}[the explicit formula, von Mangoldt 1895]\label{thm:explicit}
For $x\ge2$ and $T\ge3$,
\[
\begin{aligned}\psi_0(x)&=x-\sum_{|\gamma|<T}\frac{x^\rho}\rho-\log2\pi-\frac12\log\Bigl(1-\frac1{x^2}\Bigr)\eqbr+O\Bigl(\frac{x\log^2(xT)}T+\log x\min\Bigl(1,\frac x{T\langle x\rangle}\Bigr)\Bigr),\end{aligned}
\]
the sum running over the non-trivial zeros $\rho=\beta+i\gamma$ with $|\gamma|<T$, counted with multiplicity. The implied constant is absolute.
\end{theorem}
\begin{proof}
\emph{A good height.} Take $T'\in[T,T+1]$ as in Lemma~\ref{lem:goodT}. Proposition~\ref{prop:truncperron} with $T'$ in place of $T$ has the same error term, up to a constant. The zeros with
$T\le|\gamma|<T'$ are $O(\log T)$ in number (Corollary~\ref{cor:nunit}), each with $|\frac{x^\rho}\rho|\le\frac xT$; so we may sum over $|\gamma|<T'$ instead of $|\gamma|<T$, at a cost $O(\frac{x\log T}T)$.

\emph{The rectangle.} Let $K\ge1$ be an integer, $U=2K+1$, and let $\mathcal R$ be the rectangle with vertices $c\pm iT'$ and $-U\pm iT'$, where $c=1+\frac1{\log x}$. The function
$F(s)=-\frac{\zeta'}\zeta(s)\frac{x^s}s$ has no pole on $\partial\mathcal R$: there are no zeros on $\operatorname{Re}s=c>1$, the height $T'$ is not an ordinate, and $-U$ is odd. Inside, its poles and residues are:
\begin{itemize}
\item at $s=1$, where $-\frac{\zeta'}\zeta$ has a simple pole with residue $1$: residue $x$;
\item at each non-trivial zero $\rho$ of multiplicity $m$, where $-\frac{\zeta'}\zeta$ has residue $-m$: residue $-m\frac{x^\rho}\rho$;
\item at $s=0$, where $-\frac{\zeta'}\zeta$ is holomorphic: residue $-\frac{\zeta'}\zeta(0)=-\log2\pi$;
\item at $s=-2k$, $1\le k\le K$: residue $-\frac{x^{-2k}}{-2k}=\frac{x^{-2k}}{2k}$.
\end{itemize}
By the residue theorem, $\frac1{2\pi i}\int_{c-iT'}^{c+iT'}F$ equals the sum of these residues minus $\frac1{2\pi i}$ times the integral of $F$ over the three other sides.

\emph{The other sides.} On the left side, $|F(s)|\le C\log(U+T')\frac{x^{-U}}U$ by Lemma~\ref{lem:leftzeta}, and the side has length $2T'$: its integral tends to $0$ as $K\to\infty$. On the horizontal sides,
$|s|\ge T'$. For $-1\le\sigma\le c$, Lemma~\ref{lem:goodT} gives $|F(s)|\le C\log^2T\frac{x^\sigma}T$, and $\int_{-1}^cx^\sigma\dd\sigma\le\frac{x^c}{\log x}=\frac{ex}{\log x}$. For $\sigma\le-1$,
Lemma~\ref{lem:leftzeta} applies (the side is at height $T'\ge3$), so $|F(s)|\le C\log(|\sigma|+T+3)\frac{x^\sigma}T$, whose integral over $\sigma\le-1$ is $O(\frac{\log T}{xT})$. Altogether the horizontal sides
contribute $O\bigl(\frac{x\log^2T}{T\log x}\bigr)$.

\emph{Conclusion.} Let $K\to\infty$. The trivial zeros contribute $\sum_{k\ge1}\frac{x^{-2k}}{2k}=-\frac12\log(1-x^{-2})$. Collecting the errors, $\frac{x\log^2x}T+\frac{x\log T}T+\frac{x\log^2T}{T\log x}=O\bigl(\frac{x\log^2(xT)}T\bigr)$.
\end{proof}

Letting $T\to\infty$ with $x$ fixed, the error tends to $0$, because $\langle x\rangle>0$. This proves the formula of the introduction.

\begin{corollary}[the exact formula]\label{cor:exactexplicit}
For $x\ge2$,
\[
\begin{aligned}\psi_0(x)&=x-\lim_{T\to\infty}\sum_{|\gamma|<T}\frac{x^\rho}\rho\eqbr-\log2\pi-\frac12\log\Bigl(1-\frac1{x^2}\Bigr).\end{aligned}
\]
\end{corollary}

This answers the opening puzzle. The series $\sum\frac{x^\rho}\rho$ does not converge absolutely: $\sum\frac1{|\rho|}$ diverges (Exercise~\ref{ex:zerosums}). It converges only in the symmetric sense of
the limit, and not uniformly near a prime power, where the sum of smooth functions builds a jump, exactly like the Fourier series of a square wave. The Gibbs overshoot is
visible in Figure~\ref{fig:explicit}. Exercise~\ref{ex:cesaro} gives a smoothed version in which the series converges absolutely.

\begin{remark}[Riemann's formula]
Riemann's own formula was for $\pi(x)$, and it is more intricate. Let $J(x)=\sum_{p^k\le x}\frac1k$, which counts prime powers with weight $\frac1k$ (with the midpoint convention at the jumps). Then
$\pi(x)=\sum_{n\ge1}\frac{\mu(n)}nJ(x^{1/n})$ by Möbius inversion, and Riemann stated
\[
\begin{aligned}J(x)&=\operatorname{li}(x)-\sum_\rho\operatorname{li}(x^\rho)-\log2\eqbr+\int_x^\infty\frac{\mathrm dt}{t(t^2-1)\log t},\end{aligned}
\]
where $\operatorname{li}$ is the logarithmic integral, suitably continued to complex arguments. This version was proved by von Mangoldt too; Edwards's book (Chapters~1 and~3) tells the story and gives the
proof. The $\psi$ version is simpler because $\frac{\zeta'}\zeta$ has simple poles with integer residues, whereas $J$ corresponds to $\log\zeta$, which has logarithmic singularities.
\end{remark}

\section{The error term in the prime number theorem}

The formula says that $\psi(x)-x$ is a superposition of waves $\frac{x^\rho}\rho$, of size $\frac{x^\beta}{|\rho|}$. A zero-free region bounds every $\beta$, hence every wave. The only difficulty is that there
are infinitely many waves, and the region of Theorem~\ref{thm:zerofree} narrows as $\gamma$ grows; the truncation at height $T$ balances the two.

\begin{lemma}\label{lem:sumrho}
For $T\ge3$, $\sum_{|\gamma|<T}\frac1{|\rho|}=O(\log^2T)$.
\end{lemma}
\begin{proof}
No zero is $0$ (since $\zeta(0)\ne0$), and only finitely many have $|\gamma|<1$: they contribute $O(1)$. By Corollary~\ref{cor:nunit} and symmetry, the zeros with $k\le|\gamma|<k+1$ number $O(\log(k+2))$, and
each has $|\rho|\ge k$. The sum is $O\bigl(1+\sum_{k<T}\frac{\log(k+2)}k\bigr)=O(\log^2T)$.
\end{proof}

\begin{theorem}[de la Vallée Poussin, 1899]\label{thm:pnterror}
There is a constant $c>0$ such that
\[
\psi(x)=x+O\bigl(xe^{-c\sqrt{\log x}}\bigr).
\]
\end{theorem}
\begin{proof}
Let $x$ be large. Since $|\psi(x)-\psi_0(x)|\le\log x$ and the minimum in Theorem~\ref{thm:explicit} is at most $1$,
\[
|\psi(x)-x|\le\sum_{|\gamma|<T}\frac{x^\beta}{|\rho|}+O\Bigl(\frac{x\log^2(xT)}T+\log x\Bigr).
\]
By Theorem~\ref{thm:zerofree}, every zero with $|\gamma|<T$ has $\beta\le1-\frac{c_1}{\log T}$, so $x^\beta\le x\exp\bigl(-c_1\frac{\log x}{\log T}\bigr)$, and by Lemma~\ref{lem:sumrho}
\[
\begin{aligned}|\psi(x)-x|&=O\Bigl(x\log^2T\,e^{-c_1\frac{\log x}{\log T}}\eqbr+\frac{x\log^2(xT)}T+\log x\Bigr).\end{aligned}
\]
Now choose $T$ to balance the two main terms: $\log T=\sqrt{\log x}$. The first term becomes $x\log x\,e^{-c_1\sqrt{\log x}}$ and the second $O\bigl(x\log^2x\,e^{-\sqrt{\log x}}\bigr)$. Any $c<\min(c_1,1)$ works.
\end{proof}

The error $xe^{-c\sqrt{\log x}}$ is smaller than $\frac x{\log^kx}$ for every $k$, but larger than $x^{1-\varepsilon}$ for every $\varepsilon>0$. To get a power saving one would need a zero-free strip
$\beta\le1-\eta$ of fixed width, which nobody has.

\begin{corollary}\label{cor:pnterrorpi}
With a smaller $c>0$, $\pi(x)=\operatorname{Li}(x)+O\bigl(xe^{-c\sqrt{\log x}}\bigr)$.
\end{corollary}
\begin{proof}
By Lemma~\ref{lem:abelint} with $a_p=\log p$, $f(t)=\frac1{\log t}$ and $y=2$ (and $\theta(2)=\log2$), $\pi(x)=\frac{\theta(x)}{\log x}+\int_2^x\frac{\theta(t)}{t\log^2t}\dd t$. Integrating by parts,
$\operatorname{Li}(x)=\frac x{\log x}-\frac2{\log2}+\int_2^x\frac{\mathrm dt}{\log^2t}$. Subtracting,
\[
\begin{aligned}\pi(x)-\operatorname{Li}(x)&=\frac{\theta(x)-x}{\log x}+\frac2{\log2}\eqbr+\int_2^x\frac{\theta(t)-t}{t\log^2t}\dd t.\end{aligned}
\]
Since $\psi(x)-\theta(x)=O(\sqrt x\log x)$ (Chapter~1), $\theta(t)-t=O(te^{-c\sqrt{\log t}})$. The integral is then at most a constant times $\int_2^{\sqrt x}\mathrm dt+\int_{\sqrt x}^xe^{-c\sqrt{\log t}}\dd t\le\sqrt x+xe^{-c\sqrt{(\log x)/2}}$.
\end{proof}

\section{The Riemann hypothesis as a statement about primes}

The explicit formula makes the Riemann hypothesis concrete. If all zeros have $\beta=\frac12$, every wave has amplitude $\frac{\sqrt x}{|\rho|}$, and the waves together cannot exceed $\sqrt x$ by more than
a power of $\log x$. Conversely, a single zero with $\beta>\frac12$ would make a wave of size $x^\beta$, and nothing could cancel it forever. So the Riemann hypothesis, which is a statement about
where a complex function vanishes, is \emph{the same thing} as a statement about how regularly the primes are spread.

\begin{theorem}[von Koch, 1901]\label{thm:rhequiv}
The following are equivalent:
\begin{enumerate}[label=(\alph*)]
\item the Riemann hypothesis: every non-trivial zero has real part $\frac12$;
\item $\psi(x)=x+O(\sqrt x\log^2x)$;
\item for every $\varepsilon>0$, $\psi(x)=x+O(x^{\frac12+\varepsilon})$.
\end{enumerate}
\end{theorem}
\begin{proof}
(a)$\Rightarrow$(b). Take $T=x$ in Theorem~\ref{thm:explicit}. Each $|x^\rho|=\sqrt x$, so by Lemma~\ref{lem:sumrho} the sum over zeros is $O(\sqrt x\log^2x)$, and the error term is
$O(\frac{x\log^2(x^2)}x+\log x)=O(\log^2x)$. Also $|\psi-\psi_0|\le\log x$.

(b)$\Rightarrow$(c) is clear.

(c)$\Rightarrow$(a). For $\operatorname{Re}s>1$, Lemma~\ref{lem:abelint} with $a_n=\Lambda(n)$ and $f(t)=t^{-s}$, then $X\to\infty$ (note $\psi(X)X^{-s}\to0$ since $\psi(X)=O(X)$), gives
\[
\begin{aligned}-\frac{\zeta'}\zeta(s)&=s\int_1^\infty\frac{\psi(t)}{t^{s+1}}\dd t\\&=\frac s{s-1}+s\int_1^\infty\frac{\psi(t)-t}{t^{s+1}}\dd t.\end{aligned}
\]
Fix $\varepsilon>0$. Under (c), the last integral converges absolutely and uniformly on every half-plane $\operatorname{Re}s\ge\frac12+2\varepsilon$, so it is holomorphic on $\operatorname{Re}s>\frac12+\varepsilon$ (Tome~4,
Theorem~9.15). So the right side $G(s)$ is holomorphic on $H=\{\operatorname{Re}s>\frac12+\varepsilon\}$ except at $1$. The set $H'$ obtained from $H$ by removing $1$ and the zeros of $\zeta$ is open and connected (we
remove a discrete set), and $-\frac{\zeta'}\zeta$ and $G$ are holomorphic on $H'$ and agree for $\operatorname{Re}s>1$; by the identity theorem they agree on $H'$. If $\zeta(\rho)=0$ with $\rho\in H$, then $-\frac{\zeta'}\zeta$
has a pole at $\rho$, while $G$ is bounded near $\rho$: a contradiction. So every zero has $\beta\le\frac12+\varepsilon$, for every $\varepsilon$, hence $\beta\le\frac12$. By the symmetry $\rho\mapsto1-\overline\rho$ (Proposition~\ref{prop:xi}),
$\beta\ge\frac12$ too.
\end{proof}

\begin{think}
The proof of (c)$\Rightarrow$(a) uses (c) only to make an integral converge. Run it with the hypothesis $\psi(x)-x=O(x^\theta)$ for some $\theta<\frac12$. What do you get, and why is that absurd?
(Exercise~\ref{ex:omegahalf}.) So $\sqrt x$ is the true size of the error, up to logarithms, if the Riemann hypothesis holds; and it is the floor below which no error term can go, whether it holds or not.
\end{think}

\begin{remark}[how good is it in practice?]
Schoenfeld (1976) made (a)$\Rightarrow$(b) explicit: if the Riemann hypothesis holds, then $|\psi(x)-x|<\frac1{8\pi}\sqrt x\log^2x$ for all $x\ge73.2$. A direct computation verifies this bound up to $10^7$; the largest ratio
$\frac{|\psi(x)-x|}{\sqrt x\log^2x}$ there is $0.0365$, against $\frac1{8\pi}=0.0398$. In fact $\frac{|\psi(x)-x|}{\sqrt x}$ stays below $0.81$ for $10^3\le x\le10^7$. The truth is believed to be a little larger than
$\sqrt x$: Littlewood proved in 1914 that $\psi(x)-x$ is infinitely often larger than a constant times $\sqrt x\log\log\log x$, and infinitely often smaller than its negative.
\end{remark}

\section{Chebyshev's bias}

In 1853 Chebyshev wrote to Fuss that there seem to be more primes of the form $4n+3$ than of the form $4n+1$. Figure~\ref{fig:race} of Chapter~4 shows how right he was: up to $26860$, class $3$
is ahead or tied at every single $x$. At $x=26861$ class $1$ takes the lead for the first time, for exactly one prime; Leech found this crossing in 1957, with a computer. Up to $10^7$ class $1$ leads
at only $3406$ integers. And yet the prime number theorem for progressions says the two classes are equally large. What is going on?

The answer is a story about squares. Every odd square is $\equiv1\pmod4$. The function $\psi(x;4,1)$, which counts prime \emph{powers} with weights, is boosted by the squares of all odd primes; the
counts of primes alone are not. The natural ``fair'' object, the one whose explicit formula has no main term, is $\psi(x,\chi_{-4})=\sum_{n\le x}\chi_{-4}(n)\Lambda(n)$, and it includes the squares. So the
fair race is between prime powers, and in a race between primes alone, class $1$ has been robbed of its squares.

Let $\theta(x;q,a)=\sum_{p\le x,\ p\equiv a\,(q)}\log p$.

\begin{proposition}\label{prop:biasidentity}
For $x\ge2$, with $\chi=\chi_{-4}$,
\[
\begin{aligned}\theta(x;4,3)-\theta(x;4,1)&=\theta(\sqrt x)-\psi(x,\chi)\eqbr+O(x^{1/3}\log x).\end{aligned}
\]
In particular $\dfrac{\theta(x;4,3)-\theta(x;4,1)}{\sqrt x}=1-\dfrac{\psi(x,\chi)}{\sqrt x}+o(1)$.
\end{proposition}
\begin{proof}
Write $\psi(x,\chi)$ as $\sum_{p^k\le x}\chi(p^k)\log p$ and group the terms by the exponent $k$. For $k=1$ we get $\sum_{p\le x}\chi(p)\log p=\theta(x;4,1)-\theta(x;4,3)$. For $k=2$, $\chi(p^2)=\chi(p)^2=1$ for odd $p$ and
$0$ for $p=2$, so we get $\theta(\sqrt x)-\log2$ (when $x\ge4$). For $k\ge3$, the terms add up to at most $\sum_{3\le k\le\log_2x}\theta(x^{1/k})\le\log_2x\cdot\theta(x^{1/3})=O(x^{1/3}\log x)$ by
Proposition~\ref{prop:chebyshev}. Rearranging gives the identity. Finally $\theta(\sqrt x)\sim\sqrt x$ by the prime number theorem.
\end{proof}

So the race between primes is the race between prime powers, $-\psi(x,\chi)$, plus a head start of $\sqrt x$ for class $3$. Who wins depends on the size of $\psi(x,\chi)$ compared with $\sqrt x$.
That is exactly the scale of the waves.

\begin{remark}[the waves of $L(s,\chi_{-4})$]
The method of this chapter applies to $L(s,\chi)$ for a primitive character, with the functional equation of $L(s,\chi)$ in place of that of $\zeta$; see Davenport, Chapters~9 and~19. Since $L(s,\chi_{-4})$ has
no pole, there is no main term: for $x\ge2$,
\[
\psi_0(x,\chi_{-4})=-\lim_{T\to\infty}\sum_{|\gamma|<T}\frac{x^\rho}\rho+O(1),
\]
the sum over the non-trivial zeros of $L(s,\chi_{-4})$. The first of them, on the critical line, is $\frac12+6.0209i$. If they all lie on the line (the generalized Riemann hypothesis for $\chi_{-4}$), then,
with $x=e^u$, Proposition~\ref{prop:biasidentity} becomes
\[
\frac{\theta(e^u;4,3)-\theta(e^u;4,1)}{e^{u/2}}\approx1+\sum_\gamma\frac{e^{i\gamma u}}{\frac12+i\gamma}.
\]
The sum is an oscillation in $u$ with mean $0$, made of waves with frequencies $6.02$, $10.24$, $12.99$, \dots{} The constant $1$ is the bias. Class $1$ can lead only when the oscillation dips below $-1$,
which requires many waves to line up; it happens, but rarely.
\end{remark}

How rarely? Rubinstein and Sarnak (1994), assuming the generalized Riemann hypothesis and that the ordinates $\gamma>0$ are linearly independent over $\Q$, proved that the set of $x$ where
$\pi(x;4,3)>\pi(x;4,1)$ has logarithmic density $0.9959\ldots$: in the variable $u=\log x$, class $3$ leads $99.59\%$ of the time. Littlewood had proved in 1914, unconditionally, that the lead changes
infinitely often. Both results rest on the formula above: the waves are almost periodic in $u$, so with enough patience they all line up. The bias is there for every modulus $q$: non-squares
modulo $q$ lead squares (Exercise~\ref{ex:bias}).

\begin{think}
Primes $\equiv1\pmod3$ and $\equiv2\pmod3$ also race. Which class do you expect to lead most of the time, and why? Check with a numerical calculation, then with Exercise~\ref{ex:bias}.
\end{think}

\section{Exercises}

\begin{exercise}[$\star$]\label{ex:sumgamma}
Show that $\sum_{0<\gamma\le T}\frac1\gamma=\frac1{4\pi}\log^2T+O(\log T)$.
\end{exercise}

\begin{exercise}[$\star$]\label{ex:omegahalf}
Show that $\psi(x)-x$ is not $O(x^\theta)$ for any $\theta<\frac12$.
\end{exercise}

\begin{exercise}[$\star$]\label{ex:mertensRH}
Let $M(x)=\sum_{n\le x}\mu(n)$. Show that if $M(x)=O(x^{\frac12+\varepsilon})$ for every $\varepsilon>0$, then the Riemann hypothesis holds. (The converse is also true.)
\end{exercise}

\begin{exercise}[$\star\star$]\label{ex:shortprime}
(a) Show that for every $A>0$, the interval $\bigl]x,x+\frac x{\log^Ax}\bigr]$ contains $\sim\frac x{\log^{A+1}x}$ primes as $x\to\infty$.
(b) Show that the Riemann hypothesis implies that $]x,x+\sqrt x\log^3x]$ contains a prime for all large $x$. (Cramér, 1920.)
\end{exercise}

\begin{exercise}[$\star\star$]\label{ex:vonkoch}
Show that the Riemann hypothesis implies $\pi(x)=\operatorname{Li}(x)+O(\sqrt x\log x)$.
\end{exercise}

\begin{exercise}[$\star\star$, the signature exercise]\label{ex:bias}
(a) Show that $\theta(x;3,2)-\theta(x;3,1)=\theta(\sqrt x)-\psi(x,\chi_{-3})+O(x^{1/3}\log x)$.
(b) Modulo $8$ there are four classes of odd numbers, $1,3,5,7$, and every odd square is $\equiv1\pmod8$. Using the four characters modulo $8$, show that for $a\in\{1,3,5,7\}$,
\[
\begin{aligned}\theta(x;8,a)&=\frac{\theta(x)}4-\frac{\varepsilon_a}4\theta(\sqrt x)\\&\quad+\frac14\sum_{\chi\ne\chi_0}\chi(a)\psi(x,\chi)\eqbr+O(x^{1/3}\log x),\end{aligned}
\]
with $\varepsilon_1=3$ and $\varepsilon_a=-1$ for $a\ne1$. Which class is behind?
\end{exercise}

\begin{exercise}[$\star\star\star$]\label{ex:cesaro}
Show that for $x\ge2$,
\[
\begin{aligned}\sum_{n\le x}\Lambda(n)\Bigl(1-\frac nx\Bigr)&=\frac x2-\sum_\rho\frac{x^\rho}{\rho(\rho+1)}\eqbr-\log2\pi+\frac{\zeta'(-1)}{\zeta(-1)\,x}\\&\quad-\sum_{k\ge1}\frac{x^{-2k}}{2k(2k-1)},\end{aligned}
\]
where now the sum over zeros converges absolutely. (At $x=10$ the left side is $3.37642$, and the right side, with $200$ pairs of zeros, is $3.37645$; here $\frac{\zeta'(-1)}{\zeta(-1)}=1.98505$.)
\end{exercise}

\section{Hints}

\begin{hint}{ex:sumgamma}
Write $\frac1\gamma=\frac1T+\int_\gamma^T\frac{\mathrm dt}{t^2}$ and sum: $\sum\frac1\gamma=\frac{N(T)}T+\int\frac{N(t)}{t^2}\dd t$. Then use Theorem~\ref{thm:rvm}.
\end{hint}
\begin{hint}{ex:omegahalf}
Follow (c)$\Rightarrow$(a) in Theorem~\ref{thm:rhequiv}. Tome~4 found a zero at $\frac12+14.1347\ldots i$.
\end{hint}
\begin{hint}{ex:mertensRH}
For $\operatorname{Re}s>1$, $\frac1{\zeta(s)}$ is $\sum\mu(n)n^{-s}$, which is $s\int_1^\infty M(t)t^{-s-1}\dd t$.
\end{hint}
\begin{hint}{ex:shortprime}
Apply Theorem~\ref{thm:pnterror} (resp.\ Theorem~\ref{thm:rhequiv}) at $x+h$ and at $x$, and subtract. Remove the prime powers $p^k$, $k\ge2$, with Proposition~\ref{prop:chebyshev}.
\end{hint}
\begin{hint}{ex:vonkoch}
Use the identity for $\pi(x)-\operatorname{Li}(x)$ in the proof of Corollary~\ref{cor:pnterrorpi}.
\end{hint}
\begin{hint}{ex:bias}
(a) Copy the proof of Proposition~\ref{prop:biasidentity}; $3$ plays the role of $2$. (b) Let $\theta(x,\chi)=\sum_{p\le x}\chi(p)\log p$. Orthogonality gives $\theta(x;8,a)=\frac14\sum_\chi\chi(a)\theta(x,\chi)$, and
$\sum_{\chi\ne\chi_0}\chi(a)=4\cdot\mathbf1_{a\equiv1}-1$. The characters modulo $8$ are real.
\end{hint}
\begin{hint}{ex:cesaro}
First show that for $c>0$ the integral $\frac1{2\pi i}\int_{c-i\infty}^{c+i\infty}\frac{y^s}{s(s+1)}\dd s$ equals $(1-\frac1y)\mathbf1_{y\ge1}$. Then shift the contour as in Theorem~\ref{thm:explicit}; the factor $\frac1{s(s+1)}$ makes everything converge.
\end{hint}

\section{Solutions}

\begin{solution}{ex:sumgamma}
\idea A sum over zeros weighted by a smooth function is an integral against $\mathrm dN(t)$; the main term of $N$ does the work, its error term costs $O(\log T)$.
\writeup Finitely many ordinates lie in $]0,3[$; they contribute $O(1)$. For $3\le\gamma\le T$, $\frac1\gamma=\frac1T+\int_\gamma^T\frac{\mathrm dt}{t^2}$; summing (a finite sum),
\[
\begin{aligned}\sum_{3\le\gamma\le T}\frac1\gamma&=\frac{N(T)-N(3^-)}T\eqbr+\int_3^T\frac{N(t)-N(3^-)}{t^2}\dd t.\end{aligned}
\]
By Theorem~\ref{thm:rvm}, $\frac{N(T)}T=O(\log T)$ and $N(t)=\frac t{2\pi}\log\frac t{2\pi e}+O(\log t)$. The error contributes $\int_3^\infty\frac{\log t}{t^2}\dd t=O(1)$, and the main term
\[
\begin{aligned}\int_3^T\frac{\log(t/2\pi e)}{2\pi t}\dd t&=\frac1{4\pi}\Bigl[\log^2\frac t{2\pi e}\Bigr]_3^T\\&=\frac1{4\pi}\log^2T+O(\log T).\end{aligned}
\]
\end{solution}

\begin{solution}{ex:omegahalf}
\idea A power saving in the error term pushes the analytic continuation of $\frac{\zeta'}\zeta$ to the left, and the continuation runs into a zero.
\writeup Suppose $\psi(x)-x=O(x^\theta)$ with $\theta<\frac12$. As in the proof of Theorem~\ref{thm:rhequiv}, $\frac s{s-1}+s\int_1^\infty(\psi(t)-t)t^{-s-1}\dd t$ is holomorphic on $\operatorname{Re}s>\theta$ except at $1$, and equals
$-\frac{\zeta'}\zeta$ wherever both are holomorphic. But $\zeta$ vanishes at $\rho_1=\frac12+14.1347\ldots i$ (Tome~4 located it by a sign change of Hardy's $Z$), and $\operatorname{Re}\rho_1=\frac12>\theta$, so
$-\frac{\zeta'}\zeta$ has a pole there while the integral is bounded nearby. Contradiction. So for every $\theta<\frac12$, $|\psi(x)-x|>x^\theta$ for arbitrarily large $x$.
\end{solution}

\begin{solution}{ex:mertensRH}
\idea Same mechanism as Theorem~\ref{thm:rhequiv}, with $\frac1\zeta$ in place of $-\frac{\zeta'}\zeta$: a zero of $\zeta$ is a pole of $\frac1\zeta$.
\writeup For $\operatorname{Re}s>1$, Lemma~\ref{lem:abelint} with $a_n=\mu(n)$, $f(t)=t^{-s}$, and $M(X)X^{-s}\to0$, gives $\frac1{\zeta(s)}=s\int_1^\infty M(t)t^{-s-1}\dd t$ (Euler's product gives $\frac1\zeta=\sum\mu(n)n^{-s}$,
Tome~4, Corollary~10.4). If $M(t)=O(t^{\frac12+\varepsilon})$, the integral is holomorphic on $\operatorname{Re}s>\frac12+\varepsilon$ (Tome~4, Theorem~9.15), so $\frac1\zeta$ extends holomorphically there, by the
identity theorem on the connected set where $\zeta\ne0$. A zero of $\zeta$ with $\beta>\frac12+\varepsilon$ would be a pole of this extension. So $\beta\le\frac12$ for all zeros, and $\beta\ge\frac12$ by symmetry.
\end{solution}

\begin{solution}{ex:shortprime}
\idea The number of primes in an interval is the difference of two values of $\psi$, divided by about $\log x$. It is accurate when the length $h$ beats the error term.
\writeup (a) Let $h=\frac x{\log^Ax}$. By Theorem~\ref{thm:pnterror}, $\psi(x+h)-\psi(x)=h+O(xe^{-c\sqrt{\log x}})=h(1+o(1))$, since $e^{-c\sqrt{\log x}}\log^Ax\to0$. The prime powers $p^k\le2x$ with $k\ge2$ contribute at
most $\psi(2x)-\theta(2x)=O(\sqrt x\log x)=o(h)$, so $\theta(x+h)-\theta(x)\sim h$. Each prime in $]x,x+h]$ has $\log x<\log p\le\log2x$, so their number is between $\frac{\theta(x+h)-\theta(x)}{\log2x}$ and
$\frac{\theta(x+h)-\theta(x)}{\log x}$, both $\sim\frac h{\log x}$.

(b) Let $h=\sqrt x\log^3x$. Under the Riemann hypothesis, $\psi(x+h)-\psi(x)=h+O(\sqrt x\log^2x)$, which is positive for large $x$. Prime powers with $k\ge2$ in $]x,x+h]$: there are at most
$\frac h{2\sqrt x}+1$ squares there (consecutive squares near $x$ are $2\sqrt x$ apart), and at most one $k$-th power for each $k\ge3$ (consecutive $k$-th powers are farther than $h$ apart for large $x$).
They contribute $O\bigl((\log^3x+\log x)\log x\bigr)=o(h)$. So $\theta(x+h)-\theta(x)>0$: there is a prime.
\end{solution}

\begin{solution}{ex:vonkoch}
\idea The passage from $\psi$ to $\pi$ costs a factor $\log x$ in the main term and nothing worse in the error.
\writeup Under the Riemann hypothesis, $\theta(x)-x=\psi(x)-x+O(\sqrt x\log x)=O(\sqrt x\log^2x)$ by Theorem~\ref{thm:rhequiv} and Chapter~1. In the identity of the proof of Corollary~\ref{cor:pnterrorpi},
the first term is $O(\sqrt x\log x)$, the constant is $O(1)$, and the integral is $O\bigl(\int_2^x\frac{\sqrt t\log^2t}{t\log^2t}\dd t\bigr)=O(\sqrt x)$. So $\pi(x)-\operatorname{Li}(x)=O(\sqrt x\log x)$.
\end{solution}

\begin{solution}{ex:bias}
\idea Prime powers race fairly; primes alone are handicapped in the classes that contain squares. Modulo $8$, all odd squares fall into one class, so that class carries the whole handicap.
\writeup (a) $\chi_{-3}(p^2)=1$ for $p\ne3$ and $0$ for $p=3$, so the squares contribute $\theta(\sqrt x)-\log3$ to $\psi(x,\chi_{-3})$ (for $x\ge9$). The primes contribute $\theta(x;3,1)-\theta(x;3,2)$, and the
higher powers $O(x^{1/3}\log x)$. Rearranging gives the identity: class $2$ is favoured, by about $\sqrt x$.

(b) Let $\theta(x,\chi)=\sum_{p\le x}\chi(p)\log p$, so $\theta(x,\chi_0)=\theta(x)-\log2$. By orthogonality (Chapter~3), $\theta(x;8,a)=\frac14\sum_\chi\chi(a)\theta(x,\chi)$, the characters being real. For
$\chi\ne\chi_0$, exactly as in Proposition~\ref{prop:biasidentity}, $\psi(x,\chi)=\theta(x,\chi)+\theta(\sqrt x)-\log2+O(x^{1/3}\log x)$, since $\chi(p^2)=1$ for odd $p$. Hence
\[
\begin{aligned}\theta(x;8,a)&=\frac{\theta(x)}4+\frac14\sum_{\chi\ne\chi_0}\chi(a)\psi(x,\chi)\eqbr-\frac{\theta(\sqrt x)}4\sum_{\chi\ne\chi_0}\chi(a)+O(x^{1/3}\log x).\end{aligned}
\]
By orthogonality, $\sum_\chi\chi(a)=4\cdot\mathbf1_{a\equiv1}$, so $\sum_{\chi\ne\chi_0}\chi(a)=3$ if $a\equiv1$ and $-1$ otherwise. Class $1$ is behind each of the others by about $\sqrt x$ on average: the $\psi(x,\chi)$ are
oscillations of the same size without a preferred sign. At $x=10^6$ one finds numerically $\frac{\theta(x;8,a)-\theta(x)/4}{\sqrt x}=-0.84$, $0.34$, $-0.07$, $0.57$ for $a=1,3,5,7$: the oscillations are large. Averaged over $\log x\in[\log10^3,\log10^7]$ they
become $-0.72$, $0.23$, $0.25$, $0.23$, against the predicted $-0.75$, $0.25$, $0.25$, $0.25$.
\end{solution}

\begin{solution}{ex:cesaro}
\idea Weighting by $1-\frac nx$ instead of the sharp cut-off replaces $\frac1s$ by $\frac1{s(s+1)}$ in Perron's formula. The extra decay in $|t|$ makes the integral absolutely convergent, so no truncation
error appears, and the residue at each zero gets the extra factor $\frac1{\rho+1}$.
\writeup \emph{The kernel.} For $y>1$, shift $\int_{c-iT}^{c+iT}\frac{y^s}{s(s+1)}\dd s$ to the left as in Lemma~\ref{lem:perron}: the residues at $0$ and $-1$ are $1$ and $-\frac1y$, and the other sides are
$O(\frac{y^c}{T^2})$ and tend to $0$. For $y<1$ shift right: no residue. For $y=1$ use continuity of both sides in $y$ (dominated convergence: the integrand is bounded by $\frac{\max(y,1)^c}{|s|^2}$). So the kernel is
$(1-\frac1y)\mathbf1_{y\ge1}$, and with $y=\frac xn$ and term by term integration on $\operatorname{Re}s=c>1$,
\[
\begin{aligned}\sum_{n\le x}\Lambda(n)\Bigl(1-\frac nx\Bigr)&=\frac1{2\pi i}\int_{c-i\infty}^{c+i\infty}\Bigl(-\frac{\zeta'}\zeta(s)\Bigr)\eqbr\quad\times\frac{x^s}{s(s+1)}\dd s.\end{aligned}
\]
\emph{The shift.} Take the rectangle of Theorem~\ref{thm:explicit} with a good height $T'$. The residues are: $\frac x2$ at $1$; $-\frac{x^\rho}{\rho(\rho+1)}$ at each zero; $-\frac{\zeta'}\zeta(0)=-\log2\pi$ at $0$;
$-\frac{\zeta'}\zeta(-1)\cdot\frac{x^{-1}}{-1}=\frac{\zeta'(-1)}{\zeta(-1)x}$ at $-1$ (where $\zeta(-1)=-\frac1{12}\ne0$); and $-\frac{x^{-2k}}{(-2k)(-2k+1)}=-\frac{x^{-2k}}{2k(2k-1)}$ at $-2k$, $k\ge1$. The left side tends to $0$ as before. The horizontal sides
are now $O\bigl(\frac{x\log^2T}{T^2\log x}\bigr)$ by Lemmas~\ref{lem:goodT} and~\ref{lem:leftzeta}, and the vertical tail $|t|>T'$ is $O\bigl(\frac{x\log x}{T}\bigr)$ since $-\frac{\zeta'}\zeta(c+it)$ is bounded by $\log x+A$. Everything tends to $0$ as $T\to\infty$.
Finally $\sum_\rho\frac1{|\rho(\rho+1)|}\le\sum_\rho\frac1{|\rho|^2}<\infty$ (Exercise~\ref{ex:zerosums} gives much more), and $|x^\rho|\le x$: the series converges absolutely.
\end{solution}

\section*{Sources and further reading}

\begin{itemize}
\item H.~von Mangoldt, \emph{Zu Riemanns Abhandlung ``Über die Anzahl der Primzahlen unter einer gegebenen Grösse''}, Journal für die reine und angewandte Mathematik 114 (1895), 255--305.
\item C.-J.~de la Vallée Poussin, \emph{Sur la fonction $\zeta(s)$ de Riemann et le nombre des nombres premiers inférieurs à une limite donnée}, Mémoires couronnés de l'Académie royale de Belgique 59 (1899).
\item H.~von Koch, \emph{Sur la distribution des nombres premiers}, Acta Mathematica 24 (1901), 159--182.
\item J.~E.~Littlewood, \emph{Sur la distribution des nombres premiers}, Comptes Rendus de l'Académie des Sciences 158 (1914), 1869--1872.
\item H.~Cramér, \emph{Some theorems concerning prime numbers}, Arkiv för Matematik, Astronomi och Fysik 15 (1920), 1--33.
\item J.~Leech, \emph{Note on the distribution of prime numbers}, Journal of the London Mathematical Society 32 (1957), 56--58.
\item L.~Schoenfeld, \emph{Sharper bounds for the Chebyshev functions $\theta(x)$ and $\psi(x)$. II}, Mathematics of Computation 30 (1976), 337--360.
\item M.~Rubinstein and P.~Sarnak, \emph{Chebyshev's bias}, Experimental Mathematics 3 (1994), 173--197.
\item A.~Granville and G.~Martin, \emph{Prime number races}, American Mathematical Monthly 113 (2006), 1--33. A delightful survey, the best next read.
\item H.~Davenport, \emph{Multiplicative Number Theory}, 3rd ed., Springer GTM~74, 2000, Chapters~17--20 (the explicit formulae for $\zeta$ and $L(s,\chi)$, the prime number theorem with error term).
\item H.~L.~Montgomery and R.~C.~Vaughan, \emph{Multiplicative Number Theory I}, Cambridge, 2007, Chapters~12, 13 and~15 (explicit formulae, conditional estimates, oscillation theorems).
\item H.~M.~Edwards, \emph{Riemann's Zeta Function}, Academic Press, 1974, Chapters~1 and~10 (Riemann's formula for $\pi(x)$).
\end{itemize}
